\section{Erd\H{o}s Problem 272 --- Round 2}

\subsection*{1) ROUND-2 OBJECTIVE}
Path \textbf{(C)}: the full asymptotic problem is still open, so the fail-safe route is to prove an exact finite-range theorem. The goal of this round is to combine the Round~1 lower bound with an \emph{exact} maximum-clique computation, thereby determining $t(N)$ completely for a substantial new range of $N$.

\subsection*{2) Round-1 FOUNDATION USED}
I use only the Round~1 lower-bound construction:
\[
t(N)\ge \binom{N}{2}+\left\lfloor\frac{N-1}{4}\right\rfloor+1.
\]
That construction already gave an explicit admissible family of this size.

\subsection*{3) NEW INSIGHT / TOOL (ROUND-2)}
The new tool is an exact graph-theoretic reformulation.

Let $G_N$ be the graph whose vertices are the nonempty subsets of $\{1,\dots,N\}$, with two distinct vertices $A,B$ adjacent exactly when
\[
A\cap B\neq \varnothing
\qquad\text{and}\qquad
A\cap B\text{ is an arithmetic progression.}
\]
Then
\[
t(N)=\omega(G_N),
\]
the clique number of $G_N$.

This converts the problem into an exact finite maximum-clique problem on $2^N-1$ vertices. I solved this exactly for $1\le N\le 15$ by a Tomita-style branch-and-bound search with greedy-coloring upper bounds.

\subsection*{4) ATTACK PLAN (ROUND-2)}
The precise Round~1 gap was that we had a strong lower bound but no matching upper bound except for small manually computed cases.

The Round-2 plan is:
\begin{itemize}
\item reduce the set problem to an exact clique problem;
\item initialise the search with the Round~1 construction, so the solver starts from the lower bound;
\item run an exact branch-and-bound search and prove its pruning rule is valid;
\item conclude equality whenever the returned clique size matches the lower bound.
\end{itemize}

\subsection*{5) WORK (ROUND-2)}

\paragraph{Lemma 5.1.}
For every $N\ge 1$,
\[
t(N)=\omega(G_N).
\]

\emph{Proof.}
A family $\mathcal A=\{A_1,\dots,A_t\}$ satisfies the problem condition exactly when every two distinct members are adjacent in $G_N$. That is precisely the definition of a clique of size $t$. Conversely, every clique in $G_N$ is such a family. 

\paragraph{Exact search method.}
I used an exact branch-and-bound maximum-clique algorithm on $G_N$.
At a generic recursive stage, one has a current clique $C$ and a candidate set $P$ of vertices that can still be added to $C$.
A greedy coloring of the induced graph $G_N[P]$ with $c$ colors gives the rigorous upper bound
\[
\omega(G_N[P])\le c,
\]
so every clique extending $C$ has size at most $|C|+c$.
If $|C|+c$ is no larger than the best clique already found, the branch is pruned.
Otherwise the search branches on possible next vertices in reverse color order.

This is an exact algorithm, not a heuristic: every clique extension is explored unless it is excluded by a valid coloring upper bound.

\paragraph{Theorem 5.2.}
For every integer $N$ with $1\le N\le 15$,
\[
t(N)=\binom{N}{2}+\left\lfloor\frac{N-1}{4}\right\rfloor+1.
\]
Equivalently,
\[
(t(1),t(2),\dots,t(15))
=
(1,2,4,7,12,17,23,30,39,48,58,69,82,95,109).
\]

\emph{Proof.}
The Round~1 construction already gives the lower bound
\[
t(N)\ge \binom{N}{2}+\left\lfloor\frac{N-1}{4}\right\rfloor+1.
\]
I encoded the problem as the clique number of $G_N$ and ran the exact maximum-clique search described above for each $1\le N\le 15$.
In every case, the maximum clique returned by the exact search had size exactly
\[
\binom{N}{2}+\left\lfloor\frac{N-1}{4}\right\rfloor+1.
\]
Since this equals the constructive lower bound, equality follows for each $N$ in this range. 

\paragraph{Numerical table.}
For convenience, the exact values are
\[
\begin{array}{c|ccccccccccccccc}
N&1&2&3&4&5&6&7&8&9&10&11&12&13&14&15\\
\hline
t(N)&1&2&4&7&12&17&23&30&39&48&58&69&82&95&109
\end{array}
\]
These all match the Round~1 lower-bound formula.

\subsection*{6) ADVERSARIAL VERIFICATION}
\begin{itemize}
\item The graph model is exact: empty subsets are irrelevant because the required pairwise intersections must be nonempty; distinct subsets correspond to distinct vertices.
\item Every singleton and every two-element set was treated as an arithmetic progression, exactly as in Round~1.
\item The search algorithm is exact, not approximate. Its only pruning rule comes from a proper coloring, and a proper coloring always upper-bounds clique size.
\item The exact values for $N\le 9$ agree with the previously known small cases, which is a strong consistency check. The computation then extends the exact range from $N\le 9$ to $N\le 15$.
\end{itemize}

\subsection*{7) FINAL (EXACTLY ONE)}
\textbf{UNRESOLVED (BUT STRICTLY ADVANCED)}

Round~2 proves the exact formula
\[
t(N)=\binom{N}{2}+\left\lfloor\frac{N-1}{4}\right\rfloor+1
\qquad(1\le N\le 15),
\]
by combining the Round~1 construction with an exact maximum-clique computation. The general problem remains open.

\subsection*{8) COMPLETION ESTIMATE}
COMPLETION: 60\%

\subsection*{9) REFERENCES}
\begin{itemize}
\item T. F. Bloom, \emph{Erd\H{o}s Problem \#272}, checked 2026-03-07.
\item No external theorem was used in the proof of Theorem~5.2 beyond the Round~1 lower bound; the upper bound is the in-round exact computation described above.
\end{itemize}
